\documentclass[12pt,a4paper]{article}
        \makeatletter
        \def\input@path{{../}}
        \makeatother
        \usepackage{almeria-fichas}

        \FichaMeta{Sesion 03}{Plano por cuadrantes}{Bloque 1 · Mapa y coordenadas}{Interpretar un plano cuadriculado usando ejes y cuadrantes matematicos.}{Procedimientos}{Recordar, Aplicar, Analizar}

        \begin{document}

        \FichaCabecera

        \begin{materialesbox}
        \begin{itemize}
    \item Ficha impresa
    \item Lapiz
    \item Regla
    \item Lapis de colores
\end{itemize}
        \end{materialesbox}

        \begin{pasosbox}
        \begin{enumerate}
    \item Señala primero el eje horizontal y el vertical.
    \item Numera con cuidado antes de localizar puntos.
    \item Haz una pequena cruz en el origen.
    \item Revisa al final si has escrito primero x y despues y.
\end{enumerate}
        \end{pasosbox}

        \begin{recordatoriobox}
        \begin{itemize}
    \item El origen es el punto (0,0).
    \item El eje x es horizontal y el eje y es vertical.
    \item Cada cuadrante tiene una posicion distinta respecto al origen.
\end{itemize}
        \end{recordatoriobox}

        \begin{apoyobox}
        \begin{itemize}
    \item Marca el recorrido con el dedo: primero horizontal, despues vertical.
    \item Si te pierdes, vuelve al origen y empieza otra vez.
    \item Usa colores diferentes para ejes y cuadrantes.
\end{itemize}
        \end{apoyobox}

        \BloomSection{recordar}

\BloomExercise{recordar}{1}{Nombres basicos}{Escribe el nombre de estas partes del plano: \textbf{origen}, \textbf{eje x}, \textbf{eje y}, \textbf{cuadrante}.}{5}

\BloomExercise{recordar}{2}{Orden correcto}{Completa: en un punto cartesiano se escribe primero la coordenada \underline{\hspace{2cm}} y despues la coordenada \underline{\hspace{2cm}}.}{4}

\BloomSection{aplicar}

\BloomExercise{aplicar}{1}{Localiza puntos}{Dibuja un plano sencillo y coloca los puntos $A(2,1)$, $B(4,3)$, $C(1,4)$ y $D(3,0)$.}{7}

\BloomExercise{aplicar}{2}{Lectura de puntos}{Escribe las coordenadas de cuatro puntos que inventes en un plano de una sola zona positiva. Despues intercambia la ficha con un companero para que los sitúe.}{7}

\BloomSection{analizar}

\BloomExercise{analizar}{1}{Errores frecuentes}{Un alumno escribe $P(3,5)$ pero en el plano ha contado primero 5 hacia arriba y despues 3 a la derecha. \textbf{Analiza} el error y explica como corregirlo.}{7}

\BloomExercise{analizar}{2}{Mismo eje o mismo cuadrante}{Observa estos pares de puntos: $A(1,2)$ y $B(1,5)$; $C(2,3)$ y $D(4,3)$. Explica que tienen en comun y que cambia en cada caso.}{7}

        \begin{evidenciabox}
        \begin{itemize}
    \item Reconocer el origen, los ejes y los cuadrantes.
    \item Localizar puntos sencillos en un plano cuadriculado.
    \item Explicar como se lee un punto paso a paso.
\end{itemize}
        \end{evidenciabox}

        \end{document}
